\section{Limitations}
\label{sec:limitations}

\textbf{Limited dataset size.} MIT-BIH has 47 subjects, and after the paced records are removed the model learns from 17 people and is tested on 22. About 100\,000 beats sound like a lot, but beats from one person are near duplicates of each other, so the effective sample size is the number of people, not of beats. Every result in this report rests on a handful of patients per class.

\textbf{Class imbalance.} The training set holds 36\,112 N beats and 382 F beats, and the rare classes come from very few records (most S from 209, most F from 208). Balancing by sampling helped the V class and gave some S recall, but no balancing method can add patients. The F class in particular is not learned, and the 32 validation F beats could not guide selection.

\textbf{Differences between patients.} The inter-patient gap is the main finding: validation macro-F1 0.67, DS2 0.49, with one DS2 patient (232) accounting for most of the S errors. The model's notion of ``early'' is learned from a few rhythms and does not transfer to a slow rhythm with pauses. Beat shape, amplitude, heart rate and noise level all differ between people, and the fixed window cannot adapt to them. Conclusions from 22 test records are also fragile: moving one record would change several scores by more than the quantisation loss.

\textbf{Use of historical ECG datasets.} The recordings date from 1975 to 1979, were made in a hospital with chest electrodes on a Holter recorder, and were selected to contain interesting arrhythmias. Equipment, patient population and prevalence of each class differ from any present-day use. The annotations are expert consensus from that time and contain a small amount of label noise, especially in the noisy records that were kept.

\textbf{Absence of prospective clinical testing.} No prospective data, no clinical workflow and no clinician were involved. The model is not a medical device and nothing in this report is a claim of clinical validity. The class definitions follow the AAMI grouping, but a clinical system would need an unclassifiable output, a signal-quality output and a validated alarm policy, none of which exist here.

\textbf{Lack of hardware testing.} All deployment results are software-level: the toolchain accepted, quantised and synthesised the network, and the Python INT8 simulation agrees with the synthesiser's known-answer output. No MAX78002 board was available, so inference time, power, energy per beat, on-device accuracy and real-time behaviour are unmeasured (Section~\ref{sec:hardware-limitations}). Section~\ref{sec:hardware-validation} describes how to measure them.

\textbf{Possible domain shift to wearable ECG.} A wearable records a different lead, with dry electrodes, at a different sampling rate and resolution, with far more motion artefact, and without expert R-peak annotations. This project used the annotated R-peaks as input, so a device would additionally need a beat detector whose errors propagate into the classifier. The noise experiment used recorded noise added to hospital ECG, which is a reasonable but incomplete stand-in; electrode-motion artefact, the kind a wearable produces most, is the case where the model degrades. Resampling, re-annotation and new recordings would all be needed before the reported numbers mean anything for a wearable device.
